Finally, we prove known results on the frequentist behavior
of key quantities using our own assumptions and notation.

%%%%%%%%%%%%%%%%%%%%%
%%%%%%%%%%%%%%%%%%%%%
\begin{lem}\lemlabel{thetahat_consistent}
%
Under \assuref{bayes_clt}, $\thetahat \plim \thetatrue$, and
$\norm{\infohat^{-1}}_{op} < 2 \infoev^{-1}$ , with probability approaching one.
%
\begin{proof}
%
% For any $\epsilon >0$, let $A_\delta$ denote the event
% $\norm{\thetahat - \thetatrue}_2 > \delta$.
% By \assuref{bayes_clt} \itemref{strict_opt}, $\thetahat \in \thetaball{\delta}$
% with probability approaching one for any $\delta > 0$.

First, by strict optimality of the log likelihood (\assuref{bayes_clt}
\itemref{strict_opt}) and the boundedness of the prior density
(\assuref{bayes_clt} \itemref{prior_smooth}), we must have $\thetahat \in
\thetaball{\delta}$ with probability approaching one as $N\rightarrow \infty$,
since
%
\begin{align*}
\MoveEqLeft
\sup_{\theta \in \thetadom \setminus \thetaball{\delta}}
    \likhat(\thetatrue) - \likhat(\theta)
\\&=
\sup_{\theta \in \thetadom \setminus \thetaball{\delta}}
    \meann \ell(\x_n \vert \thetatrue) - \meann \ell(\x_n \vert \theta) +
    \frac{1}{N} \logprior(\thetatrue) - \frac{1}{N} \logprior(\theta)
\\&\ge
\sup_{\theta \in \thetadom \setminus \thetaball{\delta}}
    \meann \ell(\x_n \vert \thetatrue) - \meann \ell(\x_n \vert \theta) -
    \frac{2}{N} \sup_{\theta \in \thetadom } \abs{\logprior(\theta)}
\\&\ge
    \varepsilon +
    \frac{2}{N} \sup_{\theta \in \thetadom } \abs{\logprior(\theta)}
    & \textrm{(\assuref{bayes_clt} \itemref{strict_opt})}
\\&\ge
    \varepsilon + \ord{N^{-1}}.
    & \textrm{(\assuref{bayes_clt} \itemref{prior_proper})}
\end{align*}
%
By definition of $\thetahat$, $\likhat(\thetatrue) - \likhat(\thetahat) \le 0$,
and so $\thetahat \notin \thetadom \setminus \thetaball{\delta}$ for
sufficiently large $N$.  It therefore suffices to consider the behavior of
$\thetahat \in \thetaball{\delta}$.

By \assuref{bayes_clt} \itemref{strict_opt}, $\info$ is positive definite with
$\norm{\info^{-1}}_{op} = \norm{\likk{2}(\thetatrue)}_{op} = \infoev^{-1}$.
Choose $\delta$ small enough that $\likk{2}(\theta)$ is continuous by
\assuref{bayes_clt} \itemref{loglik_smooth}, then choose $\delta$ smaller still
so that, by continuity of the operator norm,
%
\begin{align*}
%
\sup_{\theta \in \thetaball{\delta}}
    \norm{\likk{2}(\theta)^{-1}}_{op} < \sqrt{2} \infoev^{-1}.
%
\end{align*}
%
Next, choose $\delta$ sufficiently small that, by the ULLN of
\assuref{bayes_clt} \itemref{loglik_ulln} and \lemref{ulln},
%
\begin{align*}
%
\sup_{\theta \in \thetaball{\delta}}
    \norm{\likhatk{2}(\theta)^{-1}}_{op} < 2 \infoev^{-1},
%
\end{align*}
%
with probability approaching one.  For the remainder of the proof,  we will
condition on the event that the preceding display holds.

Using \lemref{taylor_residual} to Taylor expand $\likhatk{1}$ around
$\thetahat$ and evaluate at $\thetatrue$ gives
%
\begin{align}\eqlabel{likscore_taylor_expansion}
%
\likhatk{1}(\thetatrue) ={}&
    \likhatk{1}(\thetahat) +
    \tresid{2}{\likhat, \thetahat, \thetatrue} (\thetahat - \thetatrue).
\end{align}
%

Combining the above results, the minimum eigenvalue of the
matrix $\tresid{2}{\likhat, \thetahat, \thetatrue}$ is controlled by
%
\begin{align*}
%
\inf_{v: \norm{v}_2=1}
    v^T \tresid{2}{\likhat, \thetahat, \thetatrue} v ={}&
\inf_{v: \norm{v}_2=1}
    \int_0^1 (1 - t)
        v^T \likhatk{2}(\thetahat + t (\thetatrue - \thetahat)) v dt \\
\ge&
    \inf_{v: \norm{v}_2=1}
    \inf_{\theta \in \thetaball{\delta}} v^T \likhatk{2}(\theta) v \\
\ge& \frac{\infoev}{2}.
%
\end{align*}
%
Consequently, the matrix $\tresid{2}{\likhat, \thetahat, \thetatrue}$ is invertible
with $\norm{\tresid{2}{\likhat, \thetahat, \thetatrue}^{-1}}_{op} \le
2 \infoev^{-1}$. Applying to the Taylor expansion
\eqref{likscore_taylor_expansion}, and using the fact that
$\likhatk{1}(\thetahat) = 0$, gives that
%
\begin{align*}
%
\norm{\thetahat - \thetatrue}_2 \le 4 \infoev^{-1} \norm{\likhatk{1}(\thetatrue)}_2.
%
\end{align*}
%
Finally, by a LLN applied to $\likhatk{1}$, we have
$\likhatk{1}(\thetatrue) \plim \likk{1}(\thetatrue) = 0$, concluding the proof.
%
\end{proof}
%
\end{lem}
%%%%%%%%%%%%%%%%%%%%%


\paragraph{Proof of \propref{freq_clt}.}
\begin{proof}
%
The proof uses \eqref{likscore_taylor_expansion} in the proof of
\lemref{thetahat_consistent}.  First, let $\delta_N := \norm{\thetahat - \thetatrue}_2$,
so that, by \lemref{thetahat_consistent}, $\delta_N \plim 0$.

Then, observe that
%
\begin{align*}
%
\MoveEqLeft
\norm{\tresid{2}{\likhat, \thetahat, \thetatrue} - \likk{2}(\thetatrue)}_2
\\={}&
\norm{\int_0^1 (t - 1)
\left(
    \likhatk{2}(t (\thetatrue - \thetahat) + \thetahat) dt - \likk{2}(\thetatrue)
\right) dt
}_2
\\ \le&
\sup_{\theta \in \thetaball{\delta_N}}
    \norm{\likhatk{2}(\theta) - \likk{2}(\thetatrue)}_2
\\ \le&
\sup_{\theta \in \thetaball{\delta_N}}
    \norm{\likhatk{2}(\theta) - \likhatk{2}(\thetatrue)}_2 +
\sup_{\theta \in \thetaball{\delta_N}}
    \norm{\likhatk{2}(\thetatrue) - \likk{2}(\thetatrue)}_2
\\ &\plim 0,
%
\end{align*}
%
where the limit to zero follows from continuity of $\likhatk{2}$ via
\assuref{bayes_clt} \itemref{loglik_smooth} for the first term and the ULLN via
\assuref{bayes_clt} \itemref{loglik_ulln} and \lemref{ulln} for the second term.
It follows that
%
\begin{align*}
%
\tresid{2}{\likhat, \thetahat, \thetatrue}^{-1} \plim
    \likk{2}(\thetatrue)^{-1} = -\info^{-1}.
%
\end{align*}

Next, by \assuref{bayes_clt} \itemref{loglik_ulln} and \lemref{ulln},
$\expect{\fdist(\x)}{\norm{\ellgrad{1}(\x \vert \thetatrue)}^2_2}$ is finite
with probability approaching one, so $\cov{\fdist(\x)}{\ellgrad{1}(\x \vert
\thetatrue)} = \scorecov$ is finite. By smoothness of $\lik$
(\assuref{bayes_clt} \itemref{loglik_smooth}) and the definition of
$\thetatrue$, $\expect{\fdist(\x)}{\ellgrad{1}(\x \vert \thetatrue)} = 0$. The
$\x_n$ are IID, so, by a central limit theorem,
%
\begin{align*}
%
\sqrt{N} \likhatk{1}(\thetatrue) =
    \frac{1}{\sqrt{N}} \sumn \ellgrad{1}(\x_n \vert \thetatrue)
\dlim \normdist(\cdot | 0, \scorecov).
%
\end{align*}
%
By \eqref{likscore_taylor_expansion} in the proof of
\lemref{thetahat_consistent}, with probability approaching one,
%
\begin{align*}
%
\sqrt{N}(\thetahat - \thetatrue) ={}&
    \tresid{2}{\likhat, \thetahat, \thetatrue}^{-1}
        \sqrt{N} \likhatk{1}(\thetatrue) \\
&\dlim \normdist(\cdot | 0, \info^{-1} \scorecov \info^{-1}),
%
\end{align*}
%
where the final line follows from Slutsky's theorem.
%
\end{proof}
